\slcol{अर्जुन उवाच —\\
\Index{ज्यायसी चेत्कर्मणस्ते} मता बुद्धिर्जनार्दन । \\
तत्किं कर्मणि घोरे मां नियोजयसि केशव ॥ ३.१ ॥}
\translate{arjuna uvāca —\\
jyāyasī cētkarmaṇastē matā buddhirjanārdana | \\
tatkiṁ karmaṇi ghōrē māṁ niyōjayasi kēśava ॥ 3.1 ॥}
\cquote{Arjuna said -\\
O Janardhana, you have said Buddhi Yoga is superior to ordinary actions then why are you urging me to engage especially this terrible war where all my kith and kin are present.}

\slcol{\Index{व्यामिश्रेणेव वाक्येन} बुद्धिं मोहयसीव मे ।\\
तदेकं वद निश्चित्य येन श्रेयोऽहमाप्नुयाम् ॥ ३.२ ॥}
\translate{vyāmiśrēṇēva vākyēna buddhiṁ mōhayasīva mē |\\
tadēkaṁ vada niścitya yēna śrēyōऽhamāpnuyām ॥ 3.2 ॥}
\cquote{You are confusing my mind with conflicting messages. Therefore, tell me with certainty, which one of the two paths lead me to true fulfillment.}

\newpage
\begin{mananam}{\mananamfont Mananam sloka - 1} 
\mananamtext Do I find myself in doubt about whether or not I want to be active in life? Am I half-hearted about plunging into worldly activities? Do I feel drawn to seek personal retreats, or retire away from life so as to engage deeply in spirituality? If so, is such a desire coming from a state of escapism from having to face life’s challenges and responsibilities? Is the temptation of inaction based on physical and mental lethargy?
\end{mananam}
\WritingHand\enspace{\selfReflect\small{Self Reflection}}
\begin{inspiration}{\mananamfont Inspiration}
\mananamtext True knowledge empowers us to act. Self-knowledge allows us to transcend the ego. Then, action can flow unhindered. This form of action is the most rejuvenating and shall lead to success in any field of life.
\end{inspiration}
\newpage


\slcol{श्रीभगवानुवाच —\\
\Index{लोकेऽस्मिन्द्विविधा} निष्ठा पुरा प्रोक्ता मयानघ ।\\
ज्ञानयोगेन साङ्‍ख्यानां कर्मयोगेन योगिनाम् ॥ ३.३ ॥}
\translate{śrībhagavānuvāca —\\
lōkēऽsmindvividhā niṣṭhā purā prōktā mayānagha |\\
jñānayōgēna sāṅ‍khyānāṁ karmayōgēna yōginām ॥ 3.3 ॥} 
\cquote{Sri Bhagavan said -\\
O sinless one, in the beginning of creation, I have given two kinds of paths - one is the path of knowledge for men of intellect and the other is the path of action for the yogis.}

\slcol{\Index{न कर्मणामनारम्भा}न्नैष्कर्म्यं पुरुषोऽश्नुते ।\\ 
न च संन्यसनादेव सिद्धिं समधिगच्छति ॥ ३.४ ॥}
\translate{na karmaṇāmanārambhānnaiṣkarmyaṁ puruṣōऽśnutē |\\ 
na ca saṁnyasanādēva siddhiṁ samadhigacchati ॥ 3.4 ॥}
\cquote{By not performing action, one does not attain inaction and by renouncing action, one does not achieve perfection.}

\slcol{\Index{न हि कश्चित्क्षणमपि} जातु तिष्ठत्यकर्मकृत् ।\\
कार्यते ह्यवशः कर्म सर्वः प्रकृतिजैर्गुणैः ॥ ३.५ ॥}
\translate{na hi kaścitkṣaṇamapi jātu tiṣṭhatyakarmakṛt |\\
kāryatē hyavaśaḥ karma sarvaḥ prakṛtijairguṇaiḥ ॥ 3.5 ॥}
\cquote{No one can rest for even an instant without performing action. Everyone is helplessly driven to perform actions based on their gunas, which are born of nature.}

\slcol{\Index{कर्मेन्द्रियाणि संयम्य} य आस्ते मनसा स्मरन् ।\\
इन्द्रियार्थान्विमूढात्मा मिथ्याचारः स उच्यते ॥ ३.६ ॥}
\translate{karmēndriyāṇi saṁyamya ya āstē manasā smaran |\\
indriyārthānvimūḍhātmā mithyācāraḥ sa ucyatē ॥ 3.6 ॥}
\cquote{Whoever abstains from action but allows the mind to dwell on the sensual pleasure, that deluded person is called a hypocrite.}

\newpage
\begin{mananam}{\mananamfont Mananam sloka - 3} 
\mananamtext Within the spiritual domain how can I classify my predominant personality type? Am I more extroverted and action oriented? Do I want to be of service to others? Do I feel the need to bring about changes in the world as a means of my own inner development? Or, do I tend to be introverted and contemplative? Am I ready to engage in deep study and dive within, into the quest of the Self and the Divine? If I have to balance between the two, what does that look like?
\end{mananam}
\WritingHand\enspace{\selfReflect\small{Self Reflection}}
\begin{inspiration}{\mananamfont Inspiration}
\mananamtext The broad classifications of the various spiritual paths are meant to help us carve out our own path in life. Within each of us is a variety of personalities, all of which require the right balance for progress. What is “balance” for one may not be balance for another.
\end{inspiration}
\newpage

\newpage
\begin{mananam}{\mananamfont {Mananam sloka - 5, 6}} 
\mananamtext If you are a teenager who has given up studies; a youth who is not seeking or engaged in active employment; an adult who doesn’t fulfil his or her responsibilities to the family, or otherwise shirking generative action, then you have to be aware that you are working against the natural movement and energy of nature. How then are you channelling that energy? Could your current choices with regard to the lack of action in your life be causing you mental or emotional distress? Are you lost in mental fantasies? Have you considered taking up any creative works?
\end{mananam}
\WritingHand\enspace{\selfReflect\small{Self Reflection}}
\begin{inspiration}{\mananamfont Inspiration}
\mananamtext If one is dormant and lazy it doesn’t mean that they have restrained the senses. An untrained mind cannot truly control the senses. In such cases, it is better to train the senses to engage in productive action than inaction.
\end{inspiration}
\newpage

\slcol{\Index{यस्त्विन्द्रियाणि मनसा} नियम्यारभतेऽर्जुन ।\\ 
कर्मेन्द्रियैः कर्मयोगमसक्तः स विशिष्यते ॥ ३.७ ॥}
\translate{yastvindriyāṇi manasā niyamyārabhatēऽrjuna |\\
karmēndriyaiḥ karmayōgamasaktaḥ sa viśiṣyatē ॥ 3.7 ॥}
\cquote{O Arjuna -\\
One who engages in karma yoga by disciplining the sense organs with the mind and remains unattached to the action is the one who excels.}

\slcol{\Index{नियतं कुरु कर्म} त्वं कर्म ज्यायो ह्यकर्मणः ।\\
शरीरयात्रापि च ते न प्रसिध्येदकर्मणः ॥ ३.८ ॥}
\translate{niyataṁ kuru karma tvaṁ karma jyāyō hyakarmaṇaḥ |\\ 
śarīrayātrāpi ca tē na prasidhyēdakarmaṇaḥ ॥ 3.8 ॥}
\cquote{You perform all the obligatory activities, because action is superior to inaction. And with inaction you cannot even maintain your body.}

\slcol{\Index{यज्ञार्थात्कर्मणोऽन्यत्र} लोकोऽयं कर्मबन्धनः ।\\
तदर्थं कर्म कौन्तेय मुक्तसङ्गः समाचर ॥ ३.९ ॥ }
\translate{yajñārthātkarmaṇōऽnyatra lōkōऽyaṁ karmabandhanaḥ |\\
tadarthaṁ karma kauntēya muktasaṅgaḥ samācara ॥ 3.9 ॥ }
\cquote{The world is bound by the actions, except those performed as a sacrifice. Therefore, O son of Kunti, perform action without being attached.}

\slcol{\Index{सहयज्ञाः प्रजाः} सृष्ट्वा पुरोवाच प्रजापतिः ।\\
अनेन प्रसविष्यध्वमेष वोऽस्त्विष्टकामधुक् ॥ ३.१० ॥}
\translate{sahayajñāḥ prajāḥ sṛṣṭvā purōvāca prajāpatiḥ |\\
anēna prasaviṣyadhvamēṣa vōऽstviṣṭakāmadhuk ॥ 3.10 ॥}
\cquote{In the days of yore, Prajapati created the mankind along with sacrifice(yagna) and said that "By this you can multiply and be the yielder of desired objects."}

\newpage
\begin{mananam}{\mananamfont {Mananam sloka - 7, 8}} 
\mananamtext In my line of work, how may I practice karma yoga (action without attachment to outcomes)? Can I work without motivation for reward, or do I feel a lack of energy and enthusiasm when there is no specific goal to work towards? Do I know how to act based on the demands of the moment and to break down every task to one unit at a time?
\end{mananam}
\WritingHand\enspace{\selfReflect\small{Self Reflection}}
\begin{inspiration}{\mananamfont Inspiration}
\mananamtext At the heart of karma yoga is the focus on the means rather than the end. If one gets the process right, then the end-product will take care of itself. Work for work’s sake, without a specific driving desire, is the healthiest and highest kind of action.
\end{inspiration}
\newpage

\newpage
\begin{mananam}{\mananamfont {Mananam sloka - 9}} 
\mananamtext Can I learn to value doing all my works as yagna, an act of sacrifice? Are my impulses to act triggered from a thought of result or reward? While I am doing my work, am I focussed on the moment or on the end? What about after I complete my task: Do I feel happy or sad based on the outcome?
\end{mananam}
\WritingHand\enspace{\selfReflect\small{Self Reflection}}
\begin{inspiration}{\mananamfont Inspiration}
\mananamtext The old Vedic version of yagna, often perpetuated by Vedic priests, is the offering of material things to gain merit (good karma) and heavenly rewards. Such a yagna is nothing but a trade transaction wherein one gives one thing to gain another. Karma yoga shifts this mindset toward valuing the act of giving itself.
\end{inspiration}
\newpage

\slcol{\Index{देवान्भावयतानेन ते} देवा भावयन्तु वः ।\\
परस्परं भावयन्तः श्रेयः परमवाप्स्यथ ॥ ३.११ ॥} 
\translate{dēvānbhāvayatānēna tē dēvā bhāvayantu vaḥ | \\
parasparaṁ bhāvayantaḥ śrēyaḥ paramavāpsyatha ॥ 3.11 ॥}
\cquote{Cherish the gods with this and let the gods cherish you. By cherishing each other you will attain supreme good.}

\slcol{\Index{इष्टान्भोगान्हि वो} देवा दास्यन्ते यज्ञभाविताः ।\\
तैर्दत्तानप्रदायैभ्यो यो भुङ्क्ते स्तेन एव सः ॥ ३.१२ ॥}
\translate{iṣṭānbhōgānhi vō dēvā dāsyantē yajñabhāvitāḥ |\\
tairdattānapradāyaibhyō yō bhuṅktē stēna ēva saḥ ॥ 3.12 ॥} 
\cquote{The gods, being cherished through the sacrifice (yagna) will grant desired objects. Therefore, one who enjoys those objects without offering anything in return is truly a thief.}

\slcol{\Index{यज्ञशिष्टाशिनः सन्तो} मुच्यन्ते सर्वकिल्बिषैः । \\
भुञ्जते ते त्वघं पापा ये पचन्त्यात्मकारणात् ॥ ३.१३ ॥}
\translate{yajñaśiṣṭāśinaḥ santō mucyantē sarvakilbiṣaiḥ |\\
bhuñjatē tē tvaghaṁ pāpā yē pacantyātmakāraṇāt ॥ 3.13 ॥}
\cquote{The righteous one, who partake of the remnants of the sacrifice (yagna), are freed from all sins: however, the one who cook food only for themselves incur sin.}

\slcol{\Index{अन्नाद्भवन्ति भूतानि} पर्जन्यादन्नसम्भवः ।\\
यज्ञाद्भवति पर्जन्यो यज्ञः कर्मसमुद्भवः ॥ ३.१४ ॥} 
\translate{annādbhavanti bhūtāni parjanyādannasambhavaḥ |\\
yajñādbhavati parjanyō yajñaḥ karmasamudbhavaḥ ॥ 3.14 ॥}
\cquote{From food originates the beings and from rain food originates, sacrifice (yagna) results in rain and origin of sacrifice (yagna) is action (karma).}
\newpage

\slcol{\Index{कर्म ब्रह्मोद्भवं} विद्धि ब्रह्माक्षरसमुद्भवम् ।\\
तस्मात्सर्वगतं ब्रह्म नित्यं यज्ञे प्रतिष्ठितम् ॥ ३.१५ ॥}
\translate{karma brahmōdbhavaṁ viddhi brahmākṣarasamudbhavam |\\
tasmātsarvagataṁ brahma nityaṁ yajñē pratiṣṭhitam ॥ 3.15 ॥}
\cquote{Action (karma) originates from the Vedas, and Vedas originates from the imperishable. Therefore, the all pervading Veda is eternally rooted in sacrifice (yagna).}

\slcol{\Index{एवं प्रवर्तितं चक्रं} नानुवर्तयतीह यः ।\\
अघायुरिन्द्रियारामो मोघं पार्थ स जीवति ॥ ३.१६ ॥ }
\translate{ēvaṁ pravartitaṁ cakraṁ nānuvartayatīha yaḥ |\\
aghāyurindriyārāmō mōghaṁ pārtha sa jīvati ॥ 3.16 ॥}
\cquote{One who does not follow the wheel of sacrife (yagna) and action (karma), and rejoices in sensual pleasures, lives a sinful life. O Partha, such a person lives in vain.}

\slcol{\Index{यस्त्वात्मरतिरेव} स्यादात्मतृप्तश्च मानवः । \\
आत्मन्येव च सन्तुष्टस्तस्य कार्यं न विद्यते ॥ ३.१७ ॥} 
\translate{yastvātmaratirēva syādātmatṛptaśca mānavaḥ |\\
ātmanyēva ca santuṣṭastasya kāryaṁ na vidyatē ॥ 3.17 ॥}
\cquote{But the one who is delighted in Atman, satisfied with Atman, finds content in Atman, has no obligation to perform any duties.}


\slcol{\Index{नैव तस्य कृतेनार्थो} नाकृतेनेह कश्चन ।\\
न चास्य सर्वभूतेषु कश्चिदर्थव्यपाश्रयः ॥ ३.१८ ॥} 
\translate{naiva tasya kṛtēnārthō nākṛtēnēha kaścana |\\
na cāsya sarvabhūtēṣu kaścidarthavyapāśrayaḥ ॥ 3.18 ॥}
\cquote{For one who rejoices in Atman, there is nothing to gain from action, nor from inaction – nor does such a person depend on anyone for any purpose.}

\newpage
\begin{mananam}{\mananamfont {Mananam sloka - 15}} 
\mananamtext If the universe sustains this body, is it not appropriate that I should return everything back to the universe? How can I cultivate the mindset of sacrifice and plunge into my righteous duty? Can I learn to offer everything in my work as an offering to God or Nature which provides for and nourishes my whole being?
\end{mananam}
\WritingHand\enspace{\selfReflect\small{Self Reflection}}
\begin{inspiration}{\mananamfont Inspiration}
\mananamtext The real spirit of yagna or sacrifice is to offer everything without expecting anything in return. Such an act of giving is possible when one awakens the inherent quality of shraddha (faith). Faith is what keeps the whole universe functioning and it is the essence of life. Nature nourishes those who sacrifice all, withholding nothing.  
\end{inspiration}
\newpage


\newpage
\begin{mananam}{\mananamfont {Mananam sloka - 16}} 
\mananamtext Do I detest society’s wheel of life, the rat-race and the excess focus on materialism? If so, is this because of the loser’s ‘sour grapes’ mentality, or perhaps because of laziness? If I am not participating in society’s systems, am I partaking of its benefits for myself? Am I contributing to society positively? Am I capable of swinging into action for action’s sake, without the need for selfish gratification?
\end{mananam}
\WritingHand\enspace{\selfReflect\small{Self Reflection}}
\begin{inspiration}{\mananamfont Inspiration}
\mananamtext Modern society’s focus on materialism and excessive activity is rajasic in nature. The way to balance rajas is to embrace sattva and shun tamas. The Gita points out that compared to tamas, rajas is superior, and sattva is superior to all. On the path of Karma yoga, rajas mixed with and guided by sattva is ideal.
\end{inspiration}
\newpage

\newpage
\begin{mananam}{\mananamfont {Mananam sloka - 17, 18}} 
\mananamtext Am I able to function at the relative level of this world where transactions and expectations may be natural and inevitable? If so, does it cause any internal bondage?  Am I mentally and emotionally free from any such dependence, especially from other people? What kind of personal expectations from others do I harbour?
\end{mananam}
\WritingHand\enspace{\selfReflect\small{Self Reflection}}
\begin{inspiration}{\mananamfont Inspiration}
\mananamtext A realized sage who has no expectations from anyone or anything is content with himself. Such an independent sage is freed from any obligations to society, whereas the society in which he lives has an obligation to sustain his body and his needs for the welfare of all.
\end{inspiration}
\newpage

\slcol{\Index{तस्मादसक्तः सततं} कार्यं कर्म समाचर ।\\
असक्तो ह्याचरन्कर्म परमाप्नोति पूरुषः ॥ ३.१९ ॥ }
\translate{tasmādasaktaḥ satataṁ kāryaṁ karma samācara |\\
asaktō hyācarankarma paramāpnōti pūruṣaḥ ॥ 3.19 ॥}
\cquote{Therefore, always perform obligatory duties without attachment, by performing such actions without attachment, one attains the Supreme.}

\slcol{\Index{कर्मणैव हि} संसिद्धिमास्थिता जनकादयः ।\\
लोकसङ्ग्रहमेवापि सम्पश्यन्कर्तुमर्हसि ॥ ३.२० ॥} 
\translate{karmaṇaiva hi saṁsiddhimāsthitā janakādayaḥ |\\
lōkasaṅgrahamēvāpi sampaśyankartumarhasi ॥ 3.20 ॥ }
\cquote{Janaka and others verily attained perfection through action alone, even with a view to ensure stability in the society, one should perform action.}

\slcol{\Index{यद्यदाचरति} श्रेष्ठस्तत्तदेवेतरो जनः ।\\
स यत्प्रमाणं कुरुते लोकस्तदनुवर्तते ॥ ३.२१ ॥}
\translate{yadyadācarati śrēṣṭhastattadēvētarō janaḥ |\\
sa yatpramāṇaṁ kurutē lōkastadanuvartatē ॥ 3.21 ॥}
\cquote{Whatever a superior person does, others follow; whatever standards such a person demonstrates by his action is followed by the whole world.}

\slcol{\Index{न मे पार्थास्ति} कर्तव्यं त्रिषु लोकेषु कि{\char"0E0AA}न ।\\
नानवाप्तमवाप्तव्यं वर्त एव च कर्मणि ॥ ३.२२ ॥}
\translate{na me pārthāsti kartavyaṁ triṣu lokeṣu kiñcana\\
nānavāptamavāptavyaṁ varta eva ca karmaṇi ॥ 3.22 ॥}
\cquote{O Partha, I have no duty in the three worlds, nothing that I have not attained and nothing that I have yet to attain, yet I engage in action.}

\newpage
\begin{mananam}{\mananamfont {Mananam sloka - 19}} 
\mananamtext Do my work and other actions arise from a sense of likes and dislikes? Can I see how my attachment to certain works and results is causing stress and anxiety at times? Does it lead to restlessness, sleeplessness, or other psychosomatic conditions? How can I apply the principle of immersing in work without attachment either to the kind of work or to its results?
\end{mananam}
\WritingHand\enspace{\selfReflect\small{Self Reflection}}
\begin{inspiration}{\mananamfont Inspiration}
\mananamtext Every successful person experiences the blessed state, at least momentarily, wherein the mind is one with the particular work or study at hand. Nothing and no one else remains then. The doer, the doing, and the thing to be done become one. This blessed state is the momentary experience of Nirvana, the highest joy and peace of life.
\end{inspiration}
\newpage


\newpage
\begin{mananam}{\mananamfont {Mananam sloka - 20, 21}} 
\mananamtext Who are my role models? Do they inspire me towards calm and independent action? What are my values that motivate me to work? Do they include working for the welfare of others? Can I relate to the idea of empowering myself by having a vision of the good of masses?
\end{mananam}
\WritingHand\enspace{\selfReflect\small{Self Reflection}}
\begin{inspiration}{\mananamfont Inspiration}
\mananamtext People whose spirit lives long after they are gone are those who have sacrificed their all for the welfare of the masses. Selfish desires create a lot of stress on the individual, while a desire for the masses uplifts individuals and society.
\end{inspiration}
\newpage


\newpage
\begin{mananam}{\mananamfont {Mananam sloka - 22}} 
\mananamtext Who is running this creation? What may be His or Her motive? Can I learn from the selfless functioning of every power unit in the universe – the sun, the moon, the wind, the ocean, etc.? What is it that keeps them going? What is their source and what is their nourishment? 
\end{mananam}
\WritingHand\enspace{\selfReflect\small{Self Reflection}}
\begin{inspiration}{\mananamfont Inspiration}
\mananamtext Love and compassion are the greatest energies, the twin motors that drive this universe. True love is to see everyone and everything as a part of oneself. To work for the welfare of our loved ones can never be a burden. Our life shall be in complete harmony when all our actions come from that space of love.
\end{inspiration}
\newpage

\slcol{\Index{यदि ह्यहं न} वर्तेय जातु कर्मण्यतन्द्रितः ।\\
मम वर्त्मानुवर्तन्ते मनुष्याः पार्थ सर्वशः ॥ ३.२३ ॥ }
\translate{yadi hyahaṁ na vartēya jātu karmaṇyatandritaḥ |\\
mama vartmānuvartantē manuṣyāḥ pārtha sarvaśaḥ ॥ 3.23 ॥}
\cquote{For if I did not perform action vigilantly, O Partha, everyone would follow My example in every way.}

\slcol{\Index{उत्सीदेयुरिमे लोका} न कुर्यां कर्म चेदहम् ।\\
सङ्करस्य च कर्ता स्यामुपहन्यामिमाः प्रजाः ॥ ३.२४ ॥}
\translate{utsīdēyurimē lōkā na kuryāṁ karma cēdaham |\\
saṅkarasya ca kartā syāmupahanyāmimāḥ prajāḥ ॥ 3.24 ॥}
\cquote{If I ever stop to work, these worlds will perish, I would be the cause of social disruption and be the cause of destruction of all the beings.}

\slcol{\Index{सक्ताः कर्मण्यविद्वांसो} यथा कुर्वन्ति भारत ।\\
कुर्याद्विद्वांस्तथासक्तश्चिकीर्षुर्लोकसङ्ग्रहम् ॥ ३.२५ ॥ }
\translate{saktāḥ karmaṇyavidvāṁsō yathā kurvanti bhārata |\\
kuryādvidvāṁstathāsaktaścikīrṣurlōkasaṅgraham ॥ 3.25 ॥}
\cquote{As the unenlightened act with attachment to action, O descendant of Bharata, so should the enlightened act without attachment, wishing the welfare of the world.}

\slcol{\Index{न बुद्धिभेदं} जनयेदज्ञानां कर्मसङ्गिनाम् ।\\
जोषयेत्सर्वकर्माणि विद्वान्युक्तः समाचरन् ॥ ३.२६ ॥ }
\translate{na buddhibhēdaṁ janayēdajñānāṁ karmasaṅginām |\\
jōṣayētsarvakarmāṇi vidvānyuktaḥ samācaran ॥ 3.26 ॥}
\cquote{A wise one should not confuse the ignorant ones, who are attached to action; the enlightened one, should act diligently and engage ignorant ones in all actions.}

\newpage
\begin{mananam}{\mananamfont {Mananam sloka - 25}} 
\mananamtext What is the role of knowledge in my life? How has maturity changed the way I look at life? What are the sources of wisdom in my life? Who or what can be a reliable source that can guide me to higher degrees of action? How do I shift from worldly wisdom to the wisdom of the rishis and wise ones?
\end{mananam}
\WritingHand\enspace{\selfReflect\small{Self Reflection}}
\begin{inspiration}{\mananamfont Inspiration}
\mananamtext It is instinctive to live for oneself and to fight for one’s own survival. But, being human, we are endowed with intellect. Using our wisdom we have to learn to work not just for ourselves but for all of humanity.
\end{inspiration}
\newpage

\newpage
\begin{mananam}{\mananamfont {Mananam sloka - 26}} 
\mananamtext Those who progress on the spiritual path or go to higher ashramas have a responsibility to not belittle the teachings and practices given to new aspirants. Spiritual teachings are given at different levels to different aspirants. To drop all forms of support and take refuge in the Self alone is possible only at the highest level of the spiritual ladder.
\end{mananam}
\WritingHand\enspace{\selfReflect\small{Self Reflection}}
\begin{inspiration}{\mananamfont Inspiration}
\mananamtext Focusing on the goal is sure to create stress, whereas focusing on the process empowers us to handle what is in front of us. To act without a sense of self, without likes and dislikes, is to transcend all gunas and become God-like. 
\end{inspiration}
\newpage

\slcol{\Index{प्रकृतेः क्रियमाणानि} गुणैः कर्माणि सर्वशः ।\\
अहङ्कारविमूढात्मा कर्ताहमिति मन्यते ॥ ३.२७ ॥}
\translate{prakṛtēḥ kriyamāṇāni guṇaiḥ karmāṇi sarvaśaḥ |\\ 
ahaṅkāravimūḍhātmā kartāhamiti manyatē ॥ 3.27 ॥}
\cquote{All actions are performed by the nature of qualities (gunas). One whose mind is deluded by egoism thinks, “I am the doer”.}

\slcol{\Index{तत्त्ववित्तु महाबाहो} गुणकर्मविभागयोः ।\\
गुणा गुणेषु वर्तन्त इति मत्वा न सज्जते ॥ ३.२८ ॥}
\translate{tattvavittu mahābāhō guṇakarmavibhāgayōḥ |\\
guṇā guṇēṣu vartanta iti matvā na sajjatē ॥ 3.28 ॥}
\cquote{But the one who has true insight into the realm of the gunas and their actions, understanding that the senses merely interact with the sense objects - both governed by the gunas - does not become attached.}

\slcol{\Index{प्रकृतेर्गुणसंमूढाः} सज्जन्ते गुणकर्मसु ।\\
तानकृत्स्नविदो मन्दान्कृत्स्नविन्न विचालयेत् ॥ ३.२९ ॥}
\translate{prakṛtērguṇasaṁmūḍhāḥ sajjantē guṇakarmasu |\\ 
tānakṛtsnavidō mandānkṛtsnavinna vicālayēt ॥ 3.29 ॥}
\cquote{Those who are deluded  by the gunas of Prakrithi become attached to the activities born of the gunas. But the one with true knowledge should not disturb the minds of those with dull intellect and imperfect understanding.}

\slcol{\Index{मयि सर्वाणि} कर्माणि संन्यस्याध्यात्मचेतसा ।\\
निराशीर्निर्ममो भूत्वा युध्यस्व विगतज्वरः ॥ ३.३० ॥}
\translate{mayi sarvāṇi karmāṇi saṁnyasyādhyātmacētasā |\\
nirāśīrnirmamō bhūtvā yudhyasva vigatajvaraḥ ॥ 3.30 ॥}
\cquote{Without mental agitation engage in battle and renounce all actions to Me, with mind centered on the Self, free from desire and ego.}

\newpage
\begin{mananam}{\mananamfont {Mananam sloka - 27}} 
\mananamtext Do I control the functioning of certain internal organs such as breathing and digestion, which enables the body to act? Aren’t there physiological processes supported by nature which contribute to my ability to speak? Without the physical nourishment of matter and the support of the environment, will my mind have the power to think and reason? Even those skills and abilities I claim to be my own, what has gone into their development? Considering all this, am I really the doer, the one who can claim certain achievements and accomplishments in life?
\end{mananam}
\WritingHand\enspace{\selfReflect\small{Self Reflection}}
\begin{inspiration}{\mananamfont Inspiration}
\mananamtext When we sincerely explore into the core of what constitutes our conception of “I” or “me”, we can find no fixed entity. Realization of our non-entity nature, the collapsing of the wrong notion of that what we take for granted to be our identify as the self, is the highest freedom.
\end{inspiration}
\newpage

\newpage
\begin{mananam}{\mananamfont {Mananam sloka - 30}} 
\mananamtext Whenever I put in my best efforts, do I tend to get strained? Do I get stressed and tired? Do I run into conflicts and tensions with co-workers and family members? What does “letting go” mean? Instead of compromising the intensity of my efforts, can I possibly let go of attachment to the outcome?
\end{mananam}
\WritingHand\enspace{\selfReflect\small{Self Reflection}}
\begin{inspiration}{\mananamfont Inspiration}
\mananamtext Those who have shraddha know the secret to stress-free intense effort. The moment we mentally entrust our lives to a higher force, a burden lifts from us. Then we can go about doing our best, knowing that what we deserve shall never be denied to us.
\end{inspiration}
\newpage

\slcol{\Index{ये मे मतमिदं} नित्यमनुतिष्ठन्ति मानवाः । \\
श्रद्धावन्तोऽनसूयन्तो मुच्यन्ते तेऽपि कर्मभिः ॥ ३.३१ ॥}
\translate{yē mē matamidaṁ nityamanutiṣṭhanti mānavāḥ | \\
śraddhāvantōऽnasūyantō mucyantē tēऽpi karmabhiḥ ॥ 3.31 ॥ }
\cquote{Those who follow My teachings with faith and without envy are freed from the bondage of actions.}

\slcol{\Index{ये त्वेतदभ्यसूयन्तो} नानुतिष्ठन्ति मे मतम् ।\\ 
सर्वज्ञानविमूढांस्तान्विद्धि नष्टानचेतसः ॥ ३.३२ ॥ }
\translate{yē tvētadabhyasūyantō nānutiṣṭhanti mē matam |\\
sarvajñānavimūḍhāṁstānviddhi naṣṭānacētasaḥ ॥ 3.32 ॥}
\cquote{But those who find fault with this teaching of Mine and do not follow it, devoid of knowledge and discrimination, will be deluded and doomed to destruction.}

\slcol{\Index{सदृशं चेष्टते} स्वस्याः प्रकृतेर्ज्ञानवानपि ।\\
प्रकृतिं यान्ति भूतानि निग्रहः किं करिष्यति ॥ ३.३३ ॥ }
\translate{sadṛśaṁ cēṣṭatē svasyāḥ prakṛtērjñānavānapi | \\
prakṛtiṁ yānti bhūtāni nigrahaḥ kiṁ kariṣyati ॥ 3.33 ॥ }
\cquote{Even a wise person acts according to their own nature; all beings are subject to Prakriti; What can repression accomplish?}

\slcol{\Index{इन्द्रियस्येन्द्रियस्यार्थे} रागद्वेषौ व्यवस्थितौ । \\
तयोर्न वशमागच्छेत्तौ ह्यस्य परिपन्थिनौ ॥ ३.३४ ॥}
\translate{indriyasyēndriyasyārthē rāgadvēṣau vyavasthitau |\\
tayōrna vaśamāgacchēttau hyasya paripanthinau ॥ 3.34 ॥}
\cquote{Attraction and aversion of the senses for their sense objects are natural; but one should not come under their influence, for they obtacles on the path.}

\newpage
\begin{mananam}{\mananamfont {Mananam sloka - 31, 32}} 
\mananamtext Which category do I fall into – those who have the willingness and earnestness to follow these uplifting teachings, or those who shun and lack faith in its transformational effects? Have I atleast attempted to study and understand the teachings? Have I attempted to apply them? Or am I biased against them?
\end{mananam}
\WritingHand\enspace{\selfReflect\small{Self Reflection}}
\begin{inspiration}{\mananamfont Inspiration}
\mananamtext It is a blessing to have an inherent faith in the sacred teachings, such that we are willing to conduct our life accordingly. But if such a faith is not inherent, then the Scriptures allow us to experiment with this faith and test out its validity in our own lives. Yet, if we fail to either live with inherent faith, or to test its validity, then we only suffer in having to learn the lessons of life the hard way.
\end{inspiration}
\newpage

\newpage
\begin{mananam}{\mananamfont {Mananam sloka - 33}} 
\mananamtext What are the habits which I find it hard to overcome? Am I aware of the negative side of such habits? Is there also an emotional dependence on those habits? Am I resolved to keep trying to break those habits? To what extent can I maintain mindfulness during the onslaught of temptations? How can I embrace a yogic lifestyle and those spiritual practices which work at a higher level in order to be free from these habits?
\end{mananam}
\WritingHand\enspace{\selfReflect\small{Self Reflection}}
\begin{inspiration}{\mananamfont Inspiration}
\mananamtext All wisdom and positive intentions may fail in the face of temptation. Such can be the force of addiction, that despite knowing its harmful effects, an addict is often unable to restrain himself in the presence of the object of attraction. Only by regular spiritual practices and constant mindfulness can one find permanent freedom from temptations.
\end{inspiration}
\newpage


\slcol{\Index{श्रेयान्स्वधर्मो विगुणः} परधर्मात्स्वनुष्ठितात् ।\\
स्वधर्मे निधनं श्रेयः परधर्मो भयावहः ॥ ३.३५ ॥}
\translate{śrēyānsvadharmō viguṇaḥ paradharmātsvanuṣṭhitāt |\\
svadharmē nidhanaṁ śrēyaḥ paradharmō bhayāvahaḥ ॥ 3.35 ॥}
\cquote{One’s own dharma, though devoid of merit is superior than the dharma of another performed well. It is better to die in one's own dharma; another’s dharma brings fear and danger.}

\slcol{अर्जुन उवाच —\\
\Index{अथ केन प्रयुक्तोऽयं} पापं चरति पूरुषः । \\
अनिच्छन्नपि वार्ष्णेय बलादिव नियोजितः ॥ ३.३६ ॥ }
\translate{arjuna uvāca —\\
atha kēna prayuktōऽyaṁ pāpaṁ carati pūruṣaḥ | \\
anicchannapi vārṣṇēya balādiva niyōjitaḥ ॥ 3.36 ॥}
\cquote{Arjuna said: \\
O Varshneya (Krishna), what compels one to commit sin, even unwillingly as if driven by an unknown force?}

\slcol{श्रीभगवानुवाच —\\
\Index{काम एष क्रोध} एष रजोगुणसमुद्भवः ।\\
महाशनो महापाप्मा विद्ध्येनमिह वैरिणम् ॥ ३.३७ ॥}
\translate{śrībhagavānuvāca —\\
kāma ēṣa krōdha ēṣa rajōguṇasamudbhavaḥ |\\
mahāśanō mahāpāpmā viddhyēnamiha vairiṇam ॥ 3.37 ॥}
\cquote{Bhagavan said, -\\
Desire and anger born from the rajo-guna, are the all-devouring and all-sinful; Know them to be the enemies here.}

\newpage
\begin{mananam}{\mananamfont {Mananam sloka - 35}} 
\mananamtext What are my skills and abilities in this life? Can they be successfully put to use in my line of work and activities? What traits and abilities of mine can I best utilise within my scope of work? Do I tend to aspire for others’ roles? Do I feel that I shall be more successful in some other line of work? Are those areas of work within my reach?
\end{mananam}
\WritingHand\enspace{\selfReflect\small{Self Reflection}}
\begin{inspiration}{\mananamfont Inspiration}
\mananamtext Fulfilling our rightful duties is the best means to overcome the mental conditioning of this life. These unmanifest tendencies are healthily roasted by carrying out our own true responsibilities rather than aspiring for other’s roles.
\end{inspiration}
\newpage


\newpage
\begin{mananam}{\mananamfont {Mananam sloka - 36}} 
\mananamtext Do I understand this force that impels me to act against my good intentions? Do I understand this vicious tendency in both its dormant and manifested state? Am I aware of its components through the restraint of mind and senses? Which of these two—the restraint of the mind, and that of the senses—is dominant in given situations? Am I able to restrain the force of temptations at the level of the senses? Am I able to restrain it at the level of thought?
\end{mananam}
\WritingHand\enspace{\selfReflect\small{Self Reflection}}
\begin{inspiration}{\mananamfont Inspiration}
\mananamtext Different harmful tendencies are ever lurking within us, gathered over lifetimes. It is never wise to be overconfident, thinking ourselves to be above senses impulses while giving in to temptations in a smaller, subtler degree.
\end{inspiration}
\newpage

\newpage
\begin{mananam}{\mananamfont {Mananam sloka - 37}} 
\mananamtext What are the nature of my desires and aspirations? What happens when my desires are obstructed? Do I feel anger? Do I feel fear and anxiety? Do my desires give me a sense of limitation and attachment? Are they specific to my personal happiness, or do they also include the well-being of others? Could there be something like healthy desires and unhealthy desires?
\end{mananam}
\WritingHand\enspace{\selfReflect\small{Self Reflection}}
\begin{inspiration}{\mananamfont Inspiration}
\mananamtext Our real nature is infinite. We forget that and go after objects in the world to complete us. When we feel limited then there arise selfish desires which we hope will make us feel greater and better.
\end{inspiration}
\newpage

\slcol{\Index{धूमेनाव्रियते} वह्निर्यथादर्शो मलेन च ।\\
यथोल्बेनावृतो गर्भस्तथा तेनेदमावृतम् ॥ ३.३८ ॥}
\translate{dhūmēnāvriyatē vahniryathādarśō malēna ca |\\
yathōlbēnāvṛtō garbhastathā tēnēdamāvṛtam ॥ 3.38 ॥}
\cquote{Just as fire is enveloped by smoke, a mirror is covered by dust, and an embryo is enclosed within a womb, one is shrouded by that (rajas).}


\slcol{\Index{आवृतं ज्ञानमेतेन} ज्ञानिनो नित्यवैरिणा ।\\ 
कामरूपेण कौन्तेय दुष्पूरेणानलेन च ॥ ३.३९ ॥}
\translate{āvṛtaṁ jñānamētēna jñāninō nityavairiṇā |\\
kāmarūpēṇa kauntēya duṣpūrēṇānalēna ca ॥ 3.39 ॥}
\cquote{O Kaunteya (Arjuna), wisdom is covered by eternal enemy of the wise in the form of desire, which is an unquenchable fire!}


\slcol{\Index{इन्द्रियाणि मनो} बुद्धिरस्याधिष्ठानमुच्यते ।\\
एतैर्विमोहयत्येष ज्ञानमावृत्य देहिनम् ॥ ३.४० ॥ }
\translate{indriyāṇi manō buddhirasyādhiṣṭhānamucyatē |\\
ētairvimōhayatyēṣa jñānamāvṛtya dēhinam ॥ 3.40 ॥}
\cquote{Desire, residing in the senses, mind and intellect, deludes the embodied soul by veiling its wisdom.}

\slcol{\Index{तस्मात्त्वमिन्द्रियाण्यादौ} नियम्य भरतर्षभ ।\\
पाप्मानं प्रजहिह्येनं ज्ञानविज्ञाननाशनम् ॥ ३.४१ ॥ }
\translate{tasmāttvamindriyāṇyādau niyamya bharatarṣabha |\\
pāpmānaṁ prajahihyēnaṁ jñānavijñānanāśanam ॥ 3.41 ॥}
\cquote{Therefore, O best of the Bharatas (Arjuna), first control the senses, and then destroy this sinful enemy (desire), which is the destroyer of knowledge and wisdom!}

\newpage
\begin{mananam}{\mananamfont {Mananam sloka - 39, 40}} 
\mananamtext What covers my deeper understanding and intentions? Why is it that despite my aspiration for progress, I give into temptations and laziness? Why do I lose the constant memory of working towards my higher goals?\\
Are the joys of sensual pursuits lasting? What kind of negative effect does the post-indulgent experience leave in my body and mind?
\end{mananam}
\WritingHand\enspace{\selfReflect\small{Self Reflection}}
\begin{inspiration}{\mananamfont Inspiration}
\mananamtext Wisdom at the level of intellect alone will not bring about transformation. It has to filter down to our heart and then express itself through our actions. Purity in our actions will in turn purify the core of our various tendencies.
\end{inspiration}
\newpage

\newpage
\begin{mananam}{\mananamfont {Mananam sloka - 41}} 
\mananamtext In which areas of my life do I need self-control? Am I able to restrain my senses of perception – seeing, hearing, tasting, etc? Am I able to control my senses of action – talking, walking, sexual? What is the first step I should take in exercising restraint on a daily basis?
\end{mananam}
\WritingHand\enspace{\selfReflect\small{Self Reflection}}
\begin{inspiration}{\mananamfont Inspiration}
\mananamtext Yoga is the art of restraining the unruly senses. When the senses are controlled, then wisdom is no longer enveloped. It will become available to us at the time of need to guide us towards right action.
\end{inspiration}
\newpage

\slcol{\Index{इन्द्रियाणि पराण्याहुरि}न्द्रियेभ्यः परं मनः ।\\ 
मनसस्तु परा बुद्धिर्यो बुद्धेः परतस्तु सः ॥ ३.४२ ॥ }
\translate{indriyāṇi parāṇyāhurindriyēbhyaḥ paraṁ manaḥ |\\ 
manasastu parā buddhiryō buddhēḥ paratastu saḥ ॥ 3.42 ॥ }
\cquote{They say that senses are superior, the mind is superior to the senses, the intellect is superior to the mind, and the one that is superior even to the intellect is the Atman.}

\slcol{\Index{एवं बुद्धेः परं बुद्ध्वा} संस्तभ्यात्मानमात्मना ।\\
जहि शत्रुं महाबाहो कामरूपं दुरासदम् ॥ ३.४३ ॥}
\translate{ēvaṁ buddhēḥ paraṁ buddhvā saṁstabhyātmānamātmanā |\\
jahi śatruṁ mahābāhō kāmarūpaṁ durāsadam ॥ 3.43 ॥}
\cquote{Thus, understanding that Atman is superior to the intellect and completely establishing spiritual absorption with the help of the mind, O mighty-armed Arjuna, conquer the unseizable enemy in the form of desire.}

\newpage
\begin{mananam}{\mananamfont {Mananam sloka - 42, 43}} 
\mananamtext What does it mean for me to see that the subtler is superior to and controls that which is grosser? Can I observe how my senses direct my body and its movements? Does not my mind direct my sense activity in the waking state and make the senses inactive in the sleep state when the mind itself goes to rest? \\ 
What does it mean for me to know the innermost and higher-most nature of the Self, which is superior to body, senses, mind and intellect? How can such understanding transform my outlook on life and my pursuits in the world? 
\end{mananam}
\WritingHand\enspace{\selfReflect\small{Self Reflection}}
\begin{inspiration}{\mananamfont Inspiration}
\mananamtext Our Consciousness is the most direct and highest entity in us. It is the subjective, pure, and true Self to which all other instruments (intellect, mind, senses and body) are objects. The pursuit of desires is an objective identification of this pure subject Self and desirous pursuit takes the true Self away from its real nature.
\end{inspiration}
\newpage

\begin{center}
  {\devfont ॐ  तत्सदिति  श्रीमद्भगवद्गीतासूपनिषत्सु  ब्रह्मविद्यायां योगशास्त्रे\\
  श्रीकृष्णार्जुनसंवादे  कर्मयोगो  नाम तृतीयोऽध्यायः॥\\}
  {\notosansdev\color{blue} ōṃ tatsaditi śrīmadbhagavadgītās ūpaniṣatsu \\brahmavidyāyāṃ yōgaśāstrē\\ śrīkṛṣṇārjunasaṃvādē karma-yoga\\nāma tṛtīyo'dhyāyaḥ॥}
\end{center}
\chapEnd